\documentclass[11pt]{amsart}
\synctex=1
\usepackage{amssymb, latexsym, graphicx, mathrsfs, enumerate, amsmath, amsthm, textcomp, url, tikz-cd,bbm, nccmath}
\usepackage[T1]{fontenc}
\usepackage{mathdots}
\usepackage{comment}
\usepackage{mathtools}
\usepackage{dsfont}
\usepackage[all]{xy}

\usepackage[toc,page]{appendix}
\usepackage[english]{babel}

%\usepackage{hyperref}

%\usepackage{refcheck}

%\usepackage{csquotes}

\usepackage[normalem]{ulem}

%\newcommand{\mik}[1]{{\color{blue}#1}}
%\renewcommand{\mik}[1]{{#1}}
%\newcommand{\bob}[1]{{\color{magenta}#1}}
%\newcommand{\dm}[1]{{\color{red}#1}}



\usepackage[top=2.5cm, bottom=2.5cm, left=2.6cm, right=2.6cm]{geometry}
\parskip=4pt


%\numberwithin{equation}{section}

\theoremstyle{plain}
\newtheorem{theorem}{Theorem}[section]
\newtheorem{conjecture}[theorem]{Conjecture}
\newtheorem{proposition}[theorem]{Proposition}
\newtheorem{lemma}[theorem]{Lemma}
\newtheorem{corollary}[theorem]{Corollary}
\newtheorem{question}[theorem]{Question}
\newtheorem{property}[theorem]{Property}
\newtheorem{properties}[theorem]{Properties}
\newtheorem{fact}[theorem]{Fact}

\theoremstyle{definition}

\newtheorem{example}[theorem]{Example}
\newtheorem{examples}[theorem]{Examples}
\newtheorem{remark}[theorem]{Remark}
\newtheorem{definition}[theorem]{Definition}
\newtheorem{construction}[theorem]{Construction}
\newtheorem{notation}[theorem]{Notation}
\newtheorem{subsec}[theorem]{\ha}

\newtheorem*{notation*}{Notation}

\newcommand{\SmallMatrix}[1]{\text{{%\tiny
\Small
\arraycolsep=0.4\arraycolsep\ensuremath
    {\begin{pmatrix}#1\end{pmatrix}}}}}

%Small spaces
\newcommand{\hs}{\kern 0.8pt}
\newcommand{\hssh}{\kern 1.2pt}
\newcommand{\hshs}{\kern 1.6pt}
\newcommand{\hssss}{\kern 2.0pt}
\newcommand{\ha}{{\kern 1pt}}

\newcommand{\hm}{\kern -0.8pt}
\newcommand{\hmm}{\kern -1.2pt}


\newcommand{\C}{{\mathds C}}
\newcommand{\R}{{\mathds R}}

\newcommand{\GL}{{\rm GL}}
\newcommand{\Stab}{{\rm Stab}}
\DeclareMathOperator{\Aut}{Aut}
\newcommand{\Id}{{\rm Id}}

\newcommand{\upgam}{{{}^\gamma\!}}

\newcommand{\ol}{\overline}
\newcommand{\Lie}{{\rm Lie}}
\newcommand{\PGL}{{\rm PGL}}
\newcommand{\Gr}{{\rm Gr}}
\newcommand{\Hom}{{\rm Hom}}
\newcommand{\Gm}{{\bf G}_{\rm m}}

\newcommand{\into}{{\,\hookrightarrow\,}}
\newcommand{\onto}{{\,\twoheadrightarrow\,}}
\newcommand{\isoto}{{\,\overset\sim\longrightarrow\,}}

\newcommand{\lra}{\longrightarrow}

\newcommand{\wh}{\widehat}

\newcommand{\z}{{Z}}
\renewcommand{\Lambda}{\textstyle\bigwedge}
\newcommand{\tf}{{\mathfrak t}}

\newcommand{\NN}[1]{\medskip\noindent {\bf #1.}}

%\newcommand{\C}{\mathbb{C}}
\newcommand{\Q}{\mathbb{Q}}
%\newcommand{\R}{\mathbb{R}}
\newcommand{\PP}{\mathbb{P}}
%\newcommand{\Gr}{\operatorname{Gr}}
%\newcommand{\GL}{\operatorname{GL}}
%\newcommand{\PGL}{\operatorname{PGL}}
\newcommand{\Pf}{\operatorname{Pf}}
%\newcommand{\Stab}{\operatorname{Stab}}
\newcommand{\Sing}{\operatorname{Sing}}
%\newcommand{\Aut}{\operatorname{Aut}}
\newcommand{\im}{\operatorname{im}}
\newcommand{\rk}{\operatorname{rank}}
\newcommand{\KK}{K}

\newcommand{\Mat}{{\rm Mat}}
\newcommand{\Skew}{{\rm skew}}

\newcommand{\TT}{{\sf T}}

\newcommand{\dgam}{{\hs{\gamma}}}
\newcommand{\scdot}{\!\cdot\!}




\title[Complex Lie algebra isomorphic to its complex conjugate]
{Constructing a complex Lie algebra isomorphic to its complex conjugate but not definable over reals}
\author{Mikhail Borovoi}
\address{Borovoi: Raymond and Beverly Sackler School of Mathematical Sciences,
Tel Aviv University, 6997801 Tel Aviv, Israel}
\email{borovoi@tauex.tau.ac.il}

\author{Willem A. de Graaf}
\address{De Graaf: Universit\`a di Trento, Dipartimento di Matematica, nVia Sommarive 14, 38123 Povo TN, Italy}
\email{willem.degraaf@unitn.it}

\author{Robert  M. Guralnick} %\\ {\tiny with an appendix by} D. Kerner}
\address{Guralnick: Department of Mathematics, University of Southern California, Los Angeles, CA 90089-2532, USA}
\email{guralnic@usc.edu}


\thanks{Borovoi was partially supported by the Israel Science Foundation (grant 1030/22).}
\thanks{Guralnick was partially supported by Simons Foundation Fellowship 00019819.}

\date\today

\begin{document}

\begin{abstract}
Following an idea of Demarche and using computer calculations of the authors
and ideas of LLM Claude Fable,
we construct an explicit 10-dimensional complex two-step nilpotent Lie algebra
that is isomorphic to its complex conjugate
but cannot be defined over the field of real numbers $\R$.
\end{abstract}

\begin{comment}
Abstract for arXiv:

Following an idea of Demarche and using computer calculations of the authors
and ideas of LLM Claude Fable,
we construct an explicit 10-dimensional complex two-step nilpotent Lie algebra
that is isomorphic to its complex conjugate
but cannot be defined over the field of real numbers R.

\end{comment}



\maketitle
\tableofcontents

\section{Introduction}
By the complex conjugate of a finite-dimensional complex Lie algebra $L$ we mean the complex Lie algebra obtained from $L$
by action by complex conjugation on the structure constants of $L$.
See Section \ref{s:2-C} for a construction not using a choice of a basis.
Let $L_\R$ be a be a finite-dimensional {\em real} Lie algebra; then the complex Lie algebra $L=L_\R\otimes_\R\C$
is clearly isomorphic to its complex conjugate.
Jonas Der\'e proposed the following conjecture:

\begin{conjecture}[{\cite[Conjecture 1 on page 494]{Dere}}]
\label{cj:Dere}
Assume that a complex Lie algebra $L$ is isomorphic to its complex conjugate. Then
$L$ can be defined over $\R$, that is, $L\simeq L_\R\otimes_\R\C$ for some real Lie algebra $L_\R$\hs.
\end{conjecture}

Cyril Demarche \cite[Theorem 1]{Demarche} disproved this conjecture.
He proved that {\em there exists} a $10$-dimensional complex Lie algebra $L$ is isomorphic to its complex conjugate
but cannot be defined over $\R$. However, Demarche did not produce an explicit example of such a Lie algebra $L$.
In this note we produce such an example.

\begin{theorem}\label{t:main}
Let $V_{10}=\C^{10}$ with the standard basis $e_1,\dots,e_{10}$\hs,
and let $V_{10}^*$ be the dual space with the dual basis $e_1^*,\dots e_{10}^*$.
Consider the following tensor in $(\Lambda^2 V_{10}^*)\otimes V_{10}$:

\begin{align*}
\tf=&(e_4^*\wedge e_5^*+e_3^*\wedge e_6^*)\otimes e_7\\
&+(-e_4^*\wedge e_6^*-e_3^*\wedge e_5^*+e_2^*\wedge e_6^*
+e_2^*\wedge e_5^*-e_1^*\wedge e_6^*+e_1^*\wedge e_5^*)\otimes e_8\\
&+(-e_4^*\wedge e_6^*-e_3^*\wedge e_5^*+e_2^*\wedge e_5^*
-e_2^*\wedge e_3^*-e_1^*\wedge e_6^*+e_1^*\wedge e_4^*)\otimes e_9\\
&+(-\tfrac{i}{2}\,e_2^*\wedge e_4^*+\tfrac{i}{2}\,e_1^*\wedge e_3^*
+e_1^*\wedge e_2^*)\otimes e_{10}
\end{align*}
where $i^2=-1$.
Consider the Lie bracket on $V_{10}$ induced by $\tf$:
\[ [\, ,]_\tf\colon V_{10}\times V_{10}\to V_{10}, \]
and  the obtained 10-dimensional two-step nilpotent Lie algebra $L_\tf=(V_{10},[\,,]_\tf)$\hs.
Then  (i)~$L_\tf$ is isomorphic to its complex conjugate, but (ii)~$L_\tf$ cannot be defined over $\R$.
\end{theorem}

Here (i) can be easily checked by an explicit calculation, and (ii) is  proved
in the note, following an idea of Demarche
and using computer calculations of the authors based on ideas of LLM Claude Fable.

Our $10$-dimensional Lie algebra $L=L_\tf$ is a {\em non-degenerate two-step nilpotent Lie algebra of type} $(6,4)$.
Here ``two-step nilpotent'' (another term: metabelian) means that $[\hs[L,L],L]=0$.
For details, see Section \ref{s:V,F}.

\begin{remark}
It follows from \cite[Tables 2 and 8]{GT} (see also \cite{BDG} for dimension 8 only)
that in every complex two-step nilpotent Lie algebra
of dimension $\le 8$ there exists a basis in which all structure constants are either 0 or 1,
and so the Lie algebra can be defined over $\R$ and is isomorphic to its complex conjugate.
In dimension 9, there exist complex two-step nilpotent Lie algebras $L$
such that $L$ is not isomorphic to its complex conjugate, but computer computations of Willem de Graaf
show that for complex two-step nilpotent Lie algebras of dimension 9,
Conjecture \ref{cj:Dere} holds.
Thus $d=10$ is the lowest dimension $d$ such that
for complex two-step nilpotent Lie algebras of dimension $d$,
Conjecture \ref{cj:Dere} fails.
\end{remark}

\noindent{\sc Acknowledgements.}
We thank Boris Kunyavski\u{\i} and Dmitry Kerner for helpful discussions and email correspondence.
We thank Skip Garibaldi for asking our question to LLM Claude Fable.
Claude gave an autonomous proof of Theorem \ref{t:Stab} on which the proof of Theorem \ref{t:main} is based,
but in this note we give a full proof using only the {\em ideas} of Claude, not its results.

\noindent
{\em Assumptions.} In this note, all vector spaces and Lie algebras are finite-dimensional
over a field $k$ of characteristic 0, and starting Section \ref{s:2-C} over $\C$.
By a {\em real form} of a complex Lie algebra $L$
we mean a real Lie algebra $L_\R$ such that $L_\R\otimes_\R \C\simeq L$.
We denote by $V^*$ the dual space to a vector space $V$, by $\Lambda^2 V$ the second external power of $F$,
and by $\Lambda^2 V^*\cong\big(\Lambda^2 V\big)^*$ the space of skew-symmetric bilinear forms on $V$.




\section{Constructing a two-step nilpotent Lie algebra}
\label{s:V,F}

Let $k$ be a field of characteristic 0.
Let $L$ be a finite-dimensional Lie algebra over $k$.

\begin{definition}
A Lie algebra $L$  is called {\em two-step nilpotent}
if
\[\hs [[L,L],L]=0,\]
or, which is equivalent,
\[ [L,L]\subseteq Z(L),\]
where $Z(L)$ is the center of $L$.
We say that our two-step nilpotent Lie algebra is {\em non-degenerate} if
\[ [L,L]=Z(L).\]
\end{definition}

The following two construction are due to Demarche \cite{Demarche}.

\begin{construction}\label{con:L-to-VF}
Let $L$ be a two-step nilpotent Lie algebra.
Write $Z=Z(L)$.
Then we have a canonical skew-symmetric bilinear map
\[ L\times L\to Z, \quad\ (x,y)\mapsto [x,y]\in Z,\]
which clearly induces a  skew-symmetric bilinear map
\[V_L\times V_L\to Z\quad\ \text{where $V_L=L/Z$,}\]
and a linear map
\[\psi_L\colon \Lambda^2 V_L\to Z.\]
Set
\[ F_L=\ker \psi_L\subset\Lambda^2 V_L\hs.\]
Thus from a non-degenerate two-step nilpotent Lie algebra $L$ we obtain a pair $(V_L, F_L)$.
This pair is {\em non-degenerate} in the following sense: $F_L$ does not contain a subset of the form $v\wedge V_L$
for any nonzero $v\in V_L$. Indeed, assume that $v\wedge V_L\subseteq F_L$, and let $\tilde v\in L$ be a preimage of $v$.
Then $[\tilde v,y]=0$ for all $y\in L$, whence $\tilde  v\in Z$ and $v=0$.
\end{construction}


\begin{construction}\label{con:VF-to-L}
Let $V$ be a finite-dimensional vector space over $k$,
and let $F\subset \Lambda^2 V$ be a  linear subspace
such that the pair $(V,F)$ is non-degenerate.
Consider the Lie algebra $L(V,F)$  with underlying vector space
$V\oplus (\Lambda^2 V)/F$ and with Lie bracket
\[ \big[ (v_1,x_1), (v_2,x_2)\big]=(\hs 0,\, v_1\wedge v_2+F\hs)\]
where $v_1,v_2\in V$,  $x_1,x_2\in (\Lambda^2 V)/F$, and  $v_1\wedge v_2+F\in (\Lambda^2 V)/F$.
Then for all $y_1,y_2,y_3\in L(V,F)$ we have
\begin{align}
& [y_1,y_2]=-[y_2,y_1],  \\
& [\,[y_1,y_2],y_3]=0.
\end{align}
Thus $L(V,F)$ is a two-step nilpotent Lie algebra.
We have  $Z\coloneqq \z(L(V,F))=\Lambda^2 V/F$ (because the pair $(V,F)$ is non-degenerate).
By construction, we have $[L,L]=Z$, that is, the two-step nilpotent Lie algebra $L$ is non-degenerate.
\end{construction}


\begin{proposition}
Constructions \ref{con:L-to-VF} and \ref{con:VF-to-L} induce a canonical bijection
between the set of isomorphism classes of
non-degenerate two-step nilpotent Lie algebras over $k$
and the set of isomorphism classes of non-degenerate
pairs $(V,F)$ over $k$ as in Construction \ref{con:VF-to-L}.
\end{proposition}


\begin{proof}
Clearly, these constructions induce maps on the sets of isomorphism classes.
We show that these maps are mutually inverse.

If we start with a non-degenerate pair $(V,F)$ as in Construction \ref{con:VF-to-L}, then
$L(V,F)$ is a non-degenerate two-step nilpotent Lie algebra, and clearly we have
\begin{equation}\label{e:V(L(V,F))}
(\hs V_{L(V,F)},\hs F_{L(V,F)}\hs)=(V,F).
\end{equation}
Let us now start with a non-degenerate two-step nilpotent Lie algebra $L$;
then $(V_L,F_L)$ is a non-degenerate pair.
Consider the  two-step nilpotent Lie algebra
\[ V_L\oplus \Lambda^2 V_L/F_L\hs.\]
Choose a splitting
\[\varsigma\colon V_L\to L\]
of the projection homomorphism of vector spaces
\[\pi \colon L\to L/Z=V_L\hs.\]
We set
\[\tilde V_L=\varsigma(V_L)\subset L,\quad\  \tilde F_L=(\Lambda^2\varsigma)(F_L)\subset \Lambda^2\tilde V_L\hs.\]
Then
\[ L= \tilde V_L\oplus \Lambda^2\tilde V_L/\tilde F_L\hs.\]
We have an isomorphism $V_L\isoto \tilde V_L$\,,
and the induced isomorphism $\Lambda^2 V_L\isoto \Lambda^2 \tilde V_L$
sends $F$ to $\tilde F_L$.
Thus, after choosing a splitting $\varsigma$, we obtain an isomorphism
\[ L\simeq V_L\oplus \Lambda^2 V_L/F_L\hs.\qedhere\]
\end{proof}

\begin{lemma}
Let  $(V,F)$ and $(V',F')$ be two non-degenerate pairs as in Construction \ref{con:VF-to-L},
Then  the non-degenerate two-step nilpotent Lie algebras $L(V,F)$ and $L(V',F')$ are isomorphic
if and only if $(V,F)$ and $(V',F')$ are isomorphic,
that is, there exists an isomorphism
$\varphi\colon V\isoto V'$
such that $\varphi_*(F)=F'$.
\end{lemma}

\begin{proof}
Obvious; see \eqref{e:V(L(V,F))}.
\end{proof}

\begin{corollary}\label{c:same-V}
Let $(V,F_1)$ and $(V,F_2)$ be two non-degenerate pairs with the same $V$.
Then  the non-degenerate two-step nilpotent Lie algebras $L(V,F_1)$ and $L(V,F_2)$ are isomorphic
if and only if there exists $g\in\GL(V)$
such that $g_*(F_1)=F_2$.
\end{corollary}

We consider the dual space $V^*$ to $V$. The natural pairing $V^*\times V\to k$ induces a pairing
\[ \Lambda^2 V^*\times \Lambda^2 V\to k,\]
which permits us to identify $\Lambda^2 V^* =(\Lambda^2 V)^*$.
By a $*$-pair we mean a pair  $(V, W)$ with $W\subseteq \Lambda^2 V^*$ where $V$ is a vector space over $k$.
For a pair $(V,F)$ as above, consider the orthogonal complement
\[ W=F^\bot\subseteq\Lambda^2 V^*.\]
We obtain a $*$-pair $(V, W)$.
Conversely, from a $*$-pair $(V,W)$ we obtain a pair $(V,F)$ with $F=W^\bot$.
We say that a $*$-pair $(V,W)$ in {\em non-degenerate}
if for any $x\neq 0\in V$ there exists $y\in V$ such that $\langle x\wedge y, W\rangle \neq 0$.
Then a pair $(V,F)$ is non-degenerate if and only if the corresponding $*$-pair $(V,W)$ is non-degenerate.

The group $\GL(V)$ naturally acts on $V^*$ and on $\Lambda^2 V^*$.

\begin{corollary}\label{c:same-V-*}
Let $(V,F_1)$ and $(V,F_2)$ be two non-degenerate pairs with the same $V$,
and let $(V,W_1)$ and $(V,W_2)$ be the corresponding $*$-pairs.
Then  the non-degenerate two-step nilpotent Lie algebras $L(V,F_1)$ and $L(V,F_2)$ are isomorphic
if and only if there exists $g\in\GL(V)$
such that $g^*(W_1)=W_2$.
\end{corollary}





\section{Two-step nilpotent Lie algebras over $\C$}
\label{s:2-C}

Now we take $k=\C$. Let $\gamma\in \{1,\gamma\}={\rm Gal}(\C/\R)$
denote the complex conjugation.
For a Lie algebra $L$ over $\C$, let $\dgam L$
denote the Lie algebra over $\C$ obtained by transporting structure via $\gamma$.
This means that $\dgam L$ is the same set $L$
with the same abelian group structure and the same Lie bracket,
but with the new multiplication by scalars
\[ a*x\coloneqq\upgam a\cdot x\quad\ \text{for all}\ \,a\in \C,\ x\in L.\]
For a vector space $V$ over $\C$, we define $\dgam V$ similarly.

Let $(V,F)$ be a non-degenerate pair as in Construction \ref{con:VF-to-L}.
Then
\begin{equation}\label{e:dgam}
\dgam L(V,F)\cong L(\hs\dgam V, \dgam F\hs).
\end{equation}

Let $V=\C^n=\R^n\otimes_\R\C$; then we identify $\dgam V=V$,
and $\gamma$ naturally acts on $V$.
For a $\C$-subspace $W\subset V$, write $\upgam\, W=\{\upgam x \ |\ x\in W\}$.
Then
\begin{equation}\label{e:upgam}
(\dgam V,\dgam W)\cong(V,\upgam\, W).
\end{equation}

\begin{lemma}\label{l:gamma-L}
Let $V=\C^n$ and let $(V,F)$ be a non-degenerate pair as in Construction \ref{con:VF-to-L}.
We have $\dgam L(V,F)\simeq L(V,F)$ if and only if
there exists $g\in\GL(V)$ such that $g_*(\upgam F)=F$.
\end{lemma}

\begin{proof}
By \eqref{e:dgam} and \eqref{e:upgam} we have
\[ \dgam L(V,F)\cong L(\dgam V, \dgam F)\cong L(V, \upgam F).\]
By Corollary \ref{c:same-V} we have $L(V, \upgam F)\simeq L(V,F)$
if and only if there exists $g\in\GL(V)$
such that $g_*(\upgam F)=F$, as desired.
\end{proof}

\begin{corollary}\label{c:gamma-L}
Let $V=\C^n$ and let $(V,W)$ be a non-degenerate $*$-pair. Set $F=W^\bot\subset\Lambda^2 V$.
We have $\dgam L(V,W)\simeq L(V,W)$ if and only if
there exists $g\in\GL(V)$ such that $g^*(\upgam\, W)=W$.
\end{corollary}

\begin{lemma}\label{l:real-form}
Let $V_\R=\R^n$, $V=\C^n$, and let $(V,F)$ be a non-degenerate pair as in Construction \ref{con:VF-to-L}.
The Lie algebra $L(V,F)$ admits a real form
if and only if there exists $g\in\GL(V)$
such that
\begin{equation}\label{e:exist-real}
\upgam\hs(g_*(F))= g_*(F).
\end{equation}
\end{lemma}

\begin{proof}
Assume that there exists $g\in\GL(V)$ satisfying \eqref{e:exist-real}.
Set $F'=g_*(F)$; then $\upgam F'=F'$.
Set $F'_\R=\{x\in F'\ |\ \upgam x=x\}$.
Then $F'_\R$ is a real form of $F'$, and $(V_\R,F'_\R)$ is a real form of $(V,F')$,
whence $L(V_\R, F'_\R)$ is a real form of $L(V,F')$.
By Corollary \ref{c:same-V} we have
\[ L(V,F)\cong L(V, g_*(F))=L(V,F').\]
Thus $L(V_\R, F'_\R)$ is also a real form of $L(V,F)$.

Conversely, assume that there exists a real form $L'_\R$ of $L(V,F)$;
then we have an isomorphism $L'_\R\otimes_\R\C\isoto L(V,F)$.
Set  $V'_\R=V_{L'_\R}$, $F'_\R=F_{L'_\R}$ as in Construction \ref{con:L-to-VF},
and write $V'=V'_\R\otimes _\R\C$, $F'=F'_\R\otimes _\R\C$.
Then we have isomorphisms
\[ L(V', F')\isoto L'_\R\otimes_\R\C\isoto L(V,F).\]
Choose an isomorphism $\varphi\colon V'_\R\isoto V_\R$
(existing because $V'\simeq V$, whence $\dim V'_\R=\dim V_\R$),
and set $F''_\R=\varphi_*(F'_\R)$;
then $F''_\R\subset \Lambda^2 V_\R$.
Set $F''=F''_\R\otimes_\R\C\subset (\Lambda^2 V_\R)\otimes_\R\C$; then $\upgam F''=F''$.
It follows from our constructions that
\[ L(V,F'')\simeq L(V',F')\simeq L(V,F),\]
and by Corollary \ref{c:same-V} there exists $g\in\GL(V)$ such that
$F''=g_*(F)$.
Thus
\[\upgam\hs(g_*(F))=\upgam F''=F''=g_*(F),\]
as desired.
\end{proof}

\begin{corollary}\label{c:real-form}
Let $V_\R=\R^n$, $V=\C^n$, and let $(V,W)$ be a non-degenerate $*$-pair.
The Lie algebra $L(V,W)$ admits a real form
if and only if there exists $g\in\GL(V)$
such that
\begin{equation}\label{e:exist-real-W}
\upgam\hs(g^*(W))= g^*(W).
\end{equation}
\end{corollary}







\section{A complex Lie algebra not definable over $\R$}
\label{s:F}

Let $V=\C^6$ with standard basis $e_1,\dots,e_6$, and let $i\in\C$ denote a fixed
square root of $-1$.
Then $\Lambda^2 V\simeq \C^{15}$.
Write
\[j=\SmallMatrix{0&1\\-1&0},\quad J={\rm diag}(j,j,j)\in\GL(6,\C).\]
Then $J^2=-\Id_6\subset\GL(6,\C)$ and $J_*^2=\Id_{15}\in\GL\big(\Lambda^2 V\big)$.

\begin{proposition}\label{p:real-orbit}
For $V=\C^6$, let $F\subset \Lambda^2 V$ be a subspace of codimension 4  with the following properties:
\begin{enumerate}
\item[\rm (a)] $\Stab_{\GL(V)}(F)=\{\lambda\cdot\Id_V,\ \lambda\in \C^\times\}$;
\item[\rm (b)] $J_*(\hs\upgam F)=  F$ where $\gamma\in{\rm Gal}(\C/\R)$ denotes the complex conjugation.
\end{enumerate}
Then the orbit of $F$ under $G=\GL(V)$ in the Grassmannian ${\rm Gr}(11,\Lambda^2 V)$
is stable under $\gamma$, but has no real elements,
that is, there is no $g\in\GL(V)$ such that
\begin{equation}\label{e:real}
\upgam\hs\big (g_*(F)\big)=g_*(F).
\end{equation}
\end{proposition}

\begin{proof}
It immediately follows from (b) that our orbit is $\gamma$-stable.
Assume for the sake of contradiction that there exists  $g$ as in \eqref{e:real}.
Then
\[ (\hs\upgam g\hs)_*\hs(\hs\upgam F)= g_*(F),\]
whence using (b) we obtain that
\[ \upgam g_*\hs\big(J_*^{-1} (F)\big)
=g_*(F)\]
and
\[(g^{-1}\hs \upgam g\hs J^{-1})_*\hs(F)= F.\]
Then it follows from (a) that
\[g^{-1}\hs\upgam g\hs J^{-1} = \lambda\cdot {\rm Id}_V\quad\ \text{for some}\ \, \lambda\in \C^\times\]
and
\[g^{-1}\cdot\upgam g =\lambda J.\]
Applying the operation $x\mapsto x\cdot \upgam x$ to this equality, we obtain
\[g^{-1}\hs\upgam g\, \upgam g^{-1}\hs g=\lambda\hs\upgam \lambda\hs J^2,\]
that is,
\[\Id_V=-\lambda\hs\upgam\lambda\,\Id_V\hs,\]
whence
\[\lambda\hs\upgam\lambda=-1,\]
which is a contradiction.
\end{proof}

\begin{corollary}
For $V=\C^6$ and $F\subset \Lambda^2 V$ with properties  (a) and (b)
of Proposition \ref{p:real-orbit}, let $L=L(\C^6,F)$.
Then $\upgam L\simeq L$, but $L$ cannot be defined over $\R$.
\end{corollary}

\begin{proof}
We may identify $\upgam L(\C^6,F)$ with $L(\C^6,\upgam F)$.
By Lemmas \ref{l:gamma-L} and \ref{l:real-form},
it follows from Proposition \ref{p:real-orbit}
that $\upgam L \simeq L$, but there is no real structure on $V=\C^6$
under which the subspace  $F\subset \Lambda^2 V$ is defined over $\R$,
that is, stable under complex conjugation.
Thus, our 10-dimensional two-step nilpotent Lie algebra $L=L(\C^6,F)$ of type $(6,4)$
is isomorphic to its complex conjugate, but cannot be defined over $\R$.
\end{proof}

We consider the specific 11-dimensional subspace $F\subset \Lambda^2 V=\Lambda^2\C^6$,
the orthogonal complement of the specific 4-dimensional subspace $W$ in $\Lambda^2 V^*=\Lambda^2 \C^6$
generated by the following 4 vectors:
\begin{equation}\label{e:W}
\begin{aligned}
&w_1=e_4^*\wedge e_5^*+e_3^*\wedge e_6^*\hs,\\
&w_2=-e_4^*\wedge e_6^*-e_3^*\wedge e_5^*+e_2^*\wedge e_6^*
+e_2^*\wedge e_5^*-e_1^*\wedge e_6^*+e_1^*\wedge e_5^*\hs,\\
&w_3=-e_4^*\wedge e_6^*-e_3^*\wedge e_5^*+e_2^*\wedge e_5^*
-e_2^*\wedge e_3^*-e_1^*\wedge e_6^*+e_1^*\wedge e_4^*\hs\\
&w_4=-\tfrac{i}{2}\,e_2^*\wedge e_4^*+\tfrac{i}{2}\,e_1^*\wedge e_3^*
+e_1^*\wedge e_2^*\hs.
\end{aligned}
\end{equation}


\begin{lemma}\label{l:non-deg}
The $*$-pair $(V,W)$ is non-degenerate, whence the pair $(V,F)$ is non-degenerate.
\end{lemma}

\begin{proof}
An easy computer calculation.
\end{proof}

\begin{lemma}\label{l:(b)}
We have
$J^*(\upgam\, W)= W$, and therefore $F\coloneqq W^\bot$
has property (b) of Proposition \ref{p:real-orbit}.
\end{lemma}

\begin{proof}
It is easy to check the equality in the lemma  by hand or using a computer; we have checked it both ways.
\end{proof}

\begin{theorem}\label{t:Stab}
The stabilizer of $[W]\in\Gr\big(4,\Lambda^2V^*\big)$ in $\ol G=\PGL_6(\C)$ is trivial.  That is, if
$g\in\GL_6(\C)$ satisfies $g^*(W)=W$, then $g$ is a scalar matrix.
\end{theorem}

Theorem \ref{t:Stab} will be proved in Section \ref{s:Stab-W}.

\begin{proof}[Proof of Theorem \ref{t:main} modulo Theorem \ref{t:Stab}]
By Lemma \ref{l:non-deg} the pair $(V,F)$ is non-degenerate.
It follows from Theorem \ref{t:Stab}
that the stabilizer  $\ol G_F$ in $\ol G$ of  $F=W^\bot$ is trivial,
that is, $F$ has property (a) of Proposition \ref{p:real-orbit}.
By Lemma \ref{l:(b)} the subspace  $F$ has property (b) of Proposition \ref{p:real-orbit}.
Now Theorem \ref{t:main} follows from Proposition \ref{p:real-orbit}.
\end{proof}






\section{The stabilizer of the 4-dimensional subspace}
\label{s:Stab-W}

In this section we prove Theorem \ref{t:Stab}.

Using a computer, we could easily show that the Lie algebra of the stabilizer $\ol G_W$ of $W$ in $\ol G=\PGL(6,\C)$ is trivial
(we used such computations when choosing  this specific $W$).
Therefore, $\ol G_W$ is finite.
Using the Gr\"obner basis algorithm,
we could compute the set of involutions in $\ol G_W$: no involutions.
However, a Gr\"obner basis computation of the whole finite group $\ol G_W$
or even of the set of elements of order 3 in $\ol G_W$, never ends.

Instead, we use an idea of LLM Claude Fable. Let $H=G_W$ denote the stabilizer of $W$
in $G=\GL(6,\C)$. Then $H$ naturally acts on $W$.

We identify $\Lambda^2 V^*$ with the space $\Mat_6^\Skew$ of skew-symmetric $6\times 6$ matrices.
Namely, we identify a skew-symmetric matrix $A\in \Mat_6^\Skew$ with the  skew-symmetric bilinear form
\[(u,v)\mapsto u\hs A \, v^\TT\in\C\]
where $u,v\in\C^6$ are {\em row} vectors.
Under this identification, the natural left action of $g\in\GL_6(\C)$ on $\Lambda^2V^*$ becomes
\begin{equation}\label{eq:action}
A\,\mapsto\, g A\,g^\TT.
\end{equation}
The determinant $\det(A)$ of a skew-symmetric matrix $A$ is the square of a polynomial $\Pf(A)$
called the {\em Pfaffian} of $A$; see, for instance, \cite[XV.\S9]{Lang}.
The Pfaffian satisfies the equivariance
\begin{equation}\label{eq:equivariance}
\Pf(g A g^\TT)=\det(g)\,\Pf(A).
\end{equation}
Since  $\Pf(A)^2=\det A$,  a $6\times6$ skew-symmetric matrix
$A$ has $\rk A\le4$ if and only if $\Pf(A)=0$.

In particular, for each $w=w_1,\dots,w_4$, we identify the skew-symmetric bilinear form
$w=\sum_{k<j}c_{kj}\,e_k^*\wedge e_j^*\in\Lambda^2V^*$ with the skew-symmetric
matrix $A_w=(c_{kj})\in M_6(\C)$, $A_w[k,j]=c_{kj}=-A_w[j,k]$.
For $x=(x_1,x_2,x_3,x_4)\in\C^4$, put
\[A(x)=x_1A_{w_1}+x_2A_{w_2}+x_3A_{w_3}+x_4A_{w_4}\hs.\]
Explicitly,
\begin{equation}\label{eq:Amatrix}
A(x)=\begin{pmatrix}
0 & x_4 & \tfrac{i}{2}x_4 & x_3 & x_2 & -x_2-x_3\\[1pt]
-x_4 & 0 & -x_3 & -\tfrac{i}{2}x_4 & x_2+x_3 & x_2\\[1pt]
-\tfrac{i}{2}x_4 & x_3 & 0 & 0 & -x_2-x_3 & x_1\\[1pt]
-x_3 & \tfrac{i}{2}x_4 & 0 & 0 & x_1 & -x_2-x_3\\[1pt]
-x_2 & -x_2-x_3 & x_2+x_3 & -x_1 & 0 & 0\\[1pt]
x_2+x_3 & -x_2 & -x_1 & x_2+x_3 & 0 & 0
\end{pmatrix}.
\end{equation}

\begin{lemma}\label{lem:cubic}
With $f(x):=\Pf(A(x))$ one has
\begin{equation}\label{eq:f}
f=x_1^2x_4-i\,x_1x_2x_4-2x_2^2x_3-x_2^2x_4-2x_2x_3^2-2x_2x_3x_4-x_3^2x_4
=q_3(x_1,x_2,x_3)\,x_4+q_2(x_2,x_3),
\end{equation}
where
\begin{equation}\label{eq:qc}
q_3(x_1,x_2,x_3)=x_1^2-i\,x_1x_2-x_2^2-2x_2x_3-x_3^2,\qquad
q_2(x_2,x_3)=-2x_2x_3(x_2+x_3).
\end{equation}
\end{lemma}

\begin{proof}
A Magma computation using the Magma function {\tt Pfaffian}.
\end{proof}

We consider the group $\GL(W)$, which we identify with $\GL(4,\C)$
using the basis $w_1,\dots,w_4$ of $W$.
Since by \eqref{eq:equivariance} the group $\GL(V)=\GL(6,\C)$
acting by \eqref{eq:action} preserves the Pfaffian up to a nonzero scalar,
so does the image $H|_W\subset \GL(W)$  of $H=\Stab_{\GL(V)}(W)$.
Thus
\begin{equation}\label{e:H|_W}
H|_W\subseteq\Stab_{\GL(W)}(\C^\times\scdot f) \quad\ \text{where $f$ is the cubic polynomial \eqref{eq:f}.}
\end{equation}

\begin{proposition}\label{p:Stab(f)}
$\Stab_{\GL(W)}(\C^\times\scdot f)=(\C^\times\scdot\Id_4)\cup(\C^\times\scdot \phi)$,
where
\begin{equation*}
\phi =
\begin{pmatrix}
-1 &i&0&0\\
0&1&0&0\\
0&0&1&0\\
0&0&0&1
\end{pmatrix}
\end{equation*}
\end{proposition}

\begin{proof}
A computer computation using the Gr\"obner basis algorithm.
Here, in dimension $4$, the computation does end.
\end{proof}

%A less computational alternative proof of Proposition \ref{p:Stab(f)},
%which uses the geometry of the cubic hypersurface $S\subset\PP^3$
%given by the equation $f(x)=0$, is given in Appendix \ref{app:AG}.

Now  let $x\in\C^4$, and let $A(x)$ be the skew-symmetric matrix as in \eqref{eq:Amatrix}.
Let $M=\Id_4$ or $M=\phi\in\GL(4,\C)$.
Let $G=\GL(6,\C)$.
We wish to find all $g\in G$ such that
\begin{equation}\label{e:M}
g A(x) g^\TT=A(Mx)\quad\ \text{up to scalar, for all}\ \, x\in\C^4.
\end{equation}
From this non-linear condition, following an idea of LLM Claude Fable, we pass to linear ones.

Since $A=A(x)$ and $A=A(Mx)$ are skew-symmetric,
the annihilator of $\im A$ is $\ker A$: indeed, we have
\[ n^\TT A=-(An)^\TT.\]
Note that \eqref{e:M} implies
\[\im(g\cdot A(x)\hs)=\im A(Mx),\]
because $g^\TT$ is invertible.
Hence, we have
\[ n^\TT\cdot g\cdot u=0\quad\ \text{for all}\ \, n\in \ker A(Mx),\ u\in \im A(x),\ x\in \C^4.\]
We obtain {\em linear} necessary conditions on $g$. We check them using a computer.

We obtain the following results:
for $M=\Id_4$ we have $g\in \C^\times\scdot\hs \Id_6$, and for $M=\phi$ there are no solutions.
Conversely, any $g\in \C^\times\scdot\hs \Id_6$ clearly satisfies \eqref{e:M}.
Thus we obtain:

\begin{lemma}\label{l:M}
The preimage of $(\C^\times\scdot\hs\Id_4)\cup(\C^\times\scdot \phi)\subset \GL(W)$
in $\Stab_{\GL(V)}(W)$ is $\C^\times\scdot\hs \Id_6$.
\end{lemma}

\begin{proof}[Proof of Theorem \ref{t:Stab}]
By \eqref{e:H|_W}, we have $H|_W\subseteq \Stab_{\GL(W)}(\C^\times\scdot f)$.
Using  Proposition \ref{p:Stab(f)}, we obtain that
\[H|_W\subseteq\Stab_{\GL(W)}(\C^\times\scdot f)\subseteq (\C^\times\scdot\hs\Id_4)\cup(\C^\times\scdot \phi).\]
Using Lemma \ref{l:M}, we obtain that $H\subseteq \C^\times\scdot\hs \Id_6$.
Thus $H= \C^\times\scdot\hs \Id_6$\hs, as desired.
\end{proof}

This completes the proof of  Theorem \ref{t:main}; see the end of Section \ref{s:F}.


\begin{remark}
It is known that the generic stabilizer for the action of
$\PGL_6$ on the Grassmannian $\Gr\big(4,\Lambda^2 V_6\big)$ is trivial.
This follows, for example, from \cite[Table 1.4]{GL24}.
Indeed, that result show that the stabilizer for the algebraic group is generically trivial,
and this is sufficient in characteristic zero.
In positive characteristic, one wants to show triviality of the generic
stabilizer as a group scheme.  Given the result about the algebraic group, it suffices to
check that the Lie algebra stabilizer is generically $0$.  This can be
checked at a single point  (presumably in characteristic not $2$ the
same point will work).
As in \cite{GG20}, if we are in characteristic $p$, it suffices to
check elements $x$ in the Lie algebra satisfying $x^p=x$ and $x^p=0$
and show that the dimension of the $\PGL_6$ orbit of $x$ plus the dimension of the
variety of fixed points of $x$ on the Grassmannian is strictly less
than $44$ (the dimension of the Grassmannian).
This is a straightforward linear algebra computation by considering
the various possibilities for the similarity class of $x$.

Unfortunately, this does not give rise to a Lie algebra example as in
characteristic zero.  Indeed, by the standard cohomology argument and
Lang's theorem, if $L$ is a finite-dimensional Lie algebra in characteristic $p$ with the
connected automorphism group and $L$ is isomorphic to its $q$-power
Frobenius twist, then $L$ is defined over the field of $q$ elements.
\end{remark}











\begin{comment}
\appendix
\section{Geometry of a singular cubic hypersurface in $\PP^3$}
\label{app:AG}
\smallskip

\centerline{ \em D. Kerner}
\medskip

\def\into{\stackrel{i}{\hookrightarrow}}
\def\ua{{\underline{a}}}\def\ub{{\underline{b}}}\def\uc{{\underline{c}}}\def\ul{{\underline{l}}}
\def\vC{{\vec{C}}}\def\vh{{\vec{h}}}\def\vv{{\vec{v}}}\def\vu{\vec{u}}\def\vw{\vec{w}}\def\vF{{\vec{F}}}\def\vr{{\vec{r}}}\def\vS{{\vec{S}}}
\def\vx{{\vec{x}}}\def\vy{{\vec{y}}}\def\vz{{\vec{z}}}
\newcommand{\bin}[2]{\binom{#1}{#2}}
\newcommand{\isom}{\xrightarrow[\,\smash{\raisebox{1.15ex}{\ensuremath{\scriptstyle\sim}}}\,]{}}

\newcommand{\ux}{{\underline{x}}}
\newcommand{\tx}{{\tilde x}}
\newcommand{\tA}{{\tilde A}}

Let $S=\{[x]\in\PP(W)\cong\PP^3:f(x)=0\}$ be the associated Pfaffian cubic
surface, where $\PP(W)$ carries the homogeneous coordinates $x_1,\dots,x_4$ introduced above.
In this section, using an idea of LLM Claude Fable,
we prove Proposition \ref{p:Stab(f)} using the geometry of the singular cubic $S$.

\begin{lemma}\label{lem:sing}
$\Sing(S)=\{p_0\}$ with $p_0=[0:0:0:1]=[w_4]$, and $p_0$ is an ordinary double
point \textup{(}an $A_1$ singularity\textup{)}.  Moreover, $f$ is irreducible.
\end{lemma}

\begin{proof}
The partial derivatives of $f$ are
\begin{align*}
f_{x_1}&=x_4(2x_1-ix_2),\\
f_{x_2}&=-ix_1x_4-4x_2x_3-2x_2x_4-2x_3^2-2x_3x_4,\\
f_{x_3}&=-2x_2^2-4x_2x_3-2x_2x_4-2x_3x_4,\\
f_{x_4}&=q.
\end{align*}
All four vanish at $p_0$, so $p_0\in\Sing(S)$ (note $f(p_0)=0$ since $f$ is linear
in $x_4$ with $q(0,0,0)=c(0,0)=0$).

Suppose $\nabla f=0$ at a point of $\PP^3$.  If $x_4=0$ then
$f_{x_2}=-2x_3(2x_2+x_3)=0$ and $f_{x_3}=-2x_2(x_2+2x_3)=0$, whose only common
solution is $x_2=x_3=0$ (the branches $x_3=-2x_2$, $x_2=-2x_3$ combine to force
$x_3=4x_3$); then $f_{x_4}=x_1^2=0$, so all coordinates vanish, which is not a
point of $\PP^3$.

If $x_4\ne0$, normalize $x_4=1$.  Then $f_{x_1}=0$ gives $x_1=\tfrac{i}{2}x_2$,
whence $-ix_1=\tfrac12x_2$ and
\begin{align*}
f_{x_4}&=q=-\tfrac34x_2^2-2x_2x_3-x_3^2,\\
f_{x_2}&=-\tfrac32x_2-4x_2x_3-2x_3^2-2x_3,\\
f_{x_3}&=-2\bigl(x_2^2+2x_2x_3+x_2+x_3\bigr).
\end{align*}
From $f_{x_3}=0$: if $2x_2+1=0$ then $x_2^2+2x_2x_3+x_2+x_3=\tfrac14-\tfrac12=-\tfrac14\ne0$,
impossible; hence
\begin{equation*}
x_3=-\frac{x_2(x_2+1)}{2x_2+1}.
\end{equation*}
Substituting into $q=0$ and multiplying by $(2x_2+1)^2$ yields, after
simplification,
\begin{equation*}
x_2^2\Bigl(\tfrac34(2x_2+1)^2-2(x_2+1)(2x_2+1)+(x_2+1)^2\Bigr)
=x_2^2\bigl(-x_2-\tfrac14\bigr)=0,
\end{equation*}
so $x_2=0$ or $x_2=-\tfrac14$.  If $x_2=0$ then $x_1=x_3=0$ and we recover $p_0$.
If $x_2=-\tfrac14$ then $x_3=\tfrac38$ and
$f_{x_2}=\tfrac38+\tfrac38-\tfrac{9}{32}-\tfrac34=-\tfrac{9}{32}\ne0$.
Hence $\Sing(S)=\{p_0\}$.

In the affine chart $x_4=1$ centered at $p_0$\hs, the equation reads $q+c$ with
quadratic part $q$, whose Gram matrix
\[\begin{pmatrix}1&-i/2&0\\-i/2&-1&-1\\0&-1&-1\end{pmatrix}\]
has determinant $-\tfrac14\ne0$.  Thus the projectivized tangent cone
$\{q=0\}\subset\PP^2$ is a smooth conic and $p_0$ is an $A_1$ point.  Finally, if
$f$ were reducible, $S$ would be a union of a plane and a quadric, or of three
planes, and $\Sing(S)$ would contain the (nonempty, positive-dimensional)
intersection of two components, contradicting $\Sing(S)=\{p_0\}$.
\end{proof}

\begin{remark}\label{rem:rank4}
The matrix $A(p_0)=A(0,0,0,1)$ has rank $4$ (its columns visibly span a
$4$-dimensional space).  Thus the node of $S$ does \emph{not} come from the
rank-$2$ locus $\Gr(2,V)\subset\PP(\Lambda^2V)$; rather, $\PP(W)$ is tangent to the
Pfaffian hypersurface at $w_4$.  Since the Pfaffian hypersurface is smooth away
from the rank-$2$ locus, any point of $\PP(W)$ of rank $2$ would be a singular
point of $S$; by Lemma~\ref{lem:sing} and $\rk A(p_0)=4$ there are none.  Hence
$\rk A(p)=4$ for \emph{every} $[p]\in S$.
%(The proof below only uses the exact rank verification at finitely many points; see Lemma~\ref{lem:P2}.)
\end{remark}

Lines of $\PP(W)$ through $p_0$ are parameterized by their directions
$[v]\in\PP^2$, $v=(v_1,v_2,v_3)$, via $\ell_v=\{[tv_1:tv_2:tv_3:u]\}$.

\begin{lemma}\label{lem:nodelines}
A line $\ell_v$ lies on $S$ if and only if $q(v)=c(v)=0$.  There are exactly six
such lines, with directions
\begin{equation}\label{eq:dirs}
\begin{aligned}
d_1&=(1,0,1), & d_2&=(1,0,-1), & d_3&=\Bigl(\tfrac{i+\sqrt3}{2},1,0\Bigr),\\
d_4&=\Bigl(\tfrac{i-\sqrt3}{2},1,0\Bigr), & d_5&=(0,1,-1), & d_6&=(i,1,-1).
\end{aligned}
\end{equation}
\end{lemma}

\begin{proof}
By \eqref{eq:f}, $f(tv_1,tv_2,tv_3,u)=t^2q(v)u+t^3c(v)$, which vanishes
identically in $(t,u)$ iff $q(v)=c(v)=0$.  Since $c=-2x_2x_3(x_2+x_3)$, we solve
$q=0$ on each of the three lines $\{x_2=0\}$, $\{x_3=0\}$, $\{x_2+x_3=0\}$ of
$\PP^2$: on $x_2=0$, $q=x_1^2-x_3^2$ gives $d_1,d_2$; on $x_3=0$,
$q=x_1^2-ix_1x_2-x_2^2$ gives $x_1/x_2=\tfrac{i\pm\sqrt3}{2}$, i.e.\ $d_3,d_4$;
on $x_3=-x_2$, $q=x_1(x_1-ix_2)$ gives $d_5,d_6$.  The six directions are
pairwise distinct.
\end{proof}

\begin{lemma}\label{lem:frame}
The directions $d_1,d_2,d_3,d_4$ form a projective frame of $\PP^2$: namely, writing
$\det(d_a,d_b,d_c)$ for the determinant of the matrix with rows $d_a,d_b,d_c$, we have
\begin{equation*}
\det(d_1,d_2,d_3)=2,\quad
\det(d_1,d_2,d_4)=2,\quad
\det(d_1,d_3,d_4)=\sqrt3,\quad
\det(d_2,d_3,d_4)=-\sqrt3.
\end{equation*}
\end{lemma}

\begin{proof}
Direct computation.
\end{proof}



Let $A(x)\in {\rm Mat}_{6\times 6}$ be as in \eqref{eq:Amatrix}, and take its Pfaffian
 $f(\ux):=\Pf(A(x))=x_4\cdot q_2(\ux)+q_3(\ux),$
  with $q_2(\ux)= (x_1-\frac{i}{2}x_2)^2+\frac{1}{4}x_2^2-(x_2+x_3)^2$ and $q_3(\ux)=-2x_2x_3(x_2+x_3).$

{\bf A change of variables:} we change the variable $x_1\rightsquigarrow \tx_1=x_1+\frac{i}{2}x_2.$
Accordingly, $A_2\rightsquigarrow \tA_2=A_2+\frac{i}{2} A_1.$
  In the new variables, the Pfaffian is
\begin{align*}
  f(\ux)&\coloneqq\Pf(A(x))=x_4\cdot q_2(\ux)+q_3(\ux),\quad\text{where}\\
  q_2(\ux)&=x^2_1+\frac{1}{4}x^2_2-(x_2+x_3)^2\quad\ \text{and}\\
  q_3(\ux)&= x_2x_3(x_2+x_3).
\end{align*}
Consider the cubic surface $S=V(f)\subset \PP^3.$
\begin{lemma}
Suppose an element $\phi\in \PGL_4$ preserves $S\subset \PP^3.$
Then either $\phi=\Id$, or $\phi =\Id_{1\times 1}\oplus (-\Id_{3\times 3}),$
i.e. $\phi(x_1,x_2,x_3,x_4)=(x_1,-x_2,-x_3,-x_4).$
\end{lemma}

\begin{proof}
As has been shown above, the surface $S\subset \PP^3$
has exactly one singular point, it is $[0:0:0:1].$
Thus $\phi$ must preserve this point.
Therefore, $\phi$ must preserve the set of lines on $S$ through this point.
There are exactly 6 such lines on $S,$ they are defined by
$\{q_2(\ux)=0=q_3(\ux)\}\subset \PP^3.$

Consider the intersection of these lines with the infinite plane $L_\infty:=\{x_4=0\}\subset \PP^3.$
The lines intersect $L_\infty$ at exactly 6 points.
Note that the cubic polynomial $q_3(\ux)$ splits into linear factors.
Therefore, $L_\infty\cap \{q_3(\ux)=0\}=$three lines ``at infinity".
By a direct check, these three lines (denoted $l_1,l_2,l_3$) lie on $S$  as well.

Thus the six lines through the singular point of $S$ split into three pairs,
each pair is connected at infinity by some $l_j.$
Each such pair, together with its $l_j,$ forms ``a triangle" on the surface $S.$
The surface $S$ may contain other triangles as well,
but only these three triangles contain the singular point.
Therefore, $\phi$ must permute these three triangles.

In particular, $\phi$ must preserve the infinite plane $L_\infty.$
Algebraically, we have an  equality of ideals in the ring $\phi(x_4)=(x_4)\subset \C[\ux]$.
Altogether we have for the ideals in $\C[\ux]:$
\begin{equation}
\phi(f(\ux))=(f(\ux)),\qquad \phi(x_4)=(x_4) , \qquad \phi(q_2(\ux),q_3(\ux))=(q_2(\ux),q_3(\ux)) .
\end{equation}
The latter equality of ideals corresponds to the preservation of the six lines.

The preservation of the union of lines $\cup l_j\subset L_\infty$ gives:
$\phi(q_3(\ux))=(q_3(\ux))\subset \C[x_1,x_2,x_3].$
Thus the expressions $x_2,x_3,x_2+x_3$ are permuted by $\phi$ (up to scaling).
Therefore,  $\phi(x_2,x_3)=(x_2,x_3)\subset \C[x_1,x_2,x_3].$

Finally, as $\deg(q_2)<\deg(q_3),$ the condition $\phi(q_2(\ux),q_3(\ux))=(q_2(\ux),q_3(\ux))$
implies that   $\phi(q_2(\ux))=(q_2(\ux)).$
Together with all the previous equalities of ideals, we obtain:
 \begin{equation}
\phi(x_2)=(x_2),\qquad \phi(x_3)=(x_3),\qquad \phi(x_4)=(x_4),\qquad \phi(x_1)=(x_1).
 \end{equation}
Therefore, up to scaling of $\phi$ we get:
\begin{itemize}
\item[\rm(i)] either $\phi =\Id,$ i.e. $\phi(x_j)=x_j$ ;
\item[\rm(ii)] or $\phi =\Id_{1\times 1}\oplus (-\Id_{3\times 3}),$
i.e., $\phi(x_1)= x_1 ,$ and then $\phi(x_j)= -x_j$ for $j=2,3,4.$
\end{itemize}
Returning to the original coordinates, in (ii) we obtain up to scaling
\begin{equation*}
\phi' =
\begin{pmatrix}
-1 &i&0&0\\
0&1&0&0\\
0&0&1&0\\
0&0&0&1
\end{pmatrix},
\end{equation*}
as in Proposition \ref{p:Stab(f)}.
\end{proof}

\begin{proposition}
If $g\in \PGL_6$ preserves the subspace $W\subset \Lambda^2 V,$ then $g=Id.$
\end{proposition}
In particular, the case  $\phi =\Id_{1\times 1}\oplus (-\Id_{3\times 3})$ is non-realizable.


\begin{proof}
Suppose $g$ preserves $W,$ i.e. $g\cdot {\rm Span}_\C(\{A_j\})\cdot g^\TT={\rm Span}_\C(\{A_j\}).$ Then $g A(x) g^T=A(\phi(x))$ for an element $\phi\in \PGL_4.$


Therefore, by the Lemma, either $\phi =\Id$ or $\phi =\Id_{1\times 1}\oplus (-\Id_{3\times 3}).$



 Below we restrict to the hyperplane $x_1=0$ and consider the matrix $A(x'):=A(0,x_2,x_3,x_4).$ Thus $g$ satisfies: $g A(x') g^T=-A(x').$
 Note that $\det[x_2A_2+x_3A_3+x_4A_4]\neq0.$ (E.g. because $\Pf(A(x))|_{x_1=0}\neq0.$)

 Thus we get: $g^{-\TT} A(y')^{-1} g^{-1}=-A(y')^{-1}.$ Here $y'=(y_2,y_3,y_4)$ is a new (formal) variable and
 the entries of $A(y')^{-1}$ lie in the field $\C(y).$
 Together with $g A(x') g^\TT=-A(x')$ we get:

\begin{equation}
g\cdot A(x')A(y')^{-1}=A(x')A(y')^{-1}\cdot g.
\end{equation}
This condition holds for all the points $x'\in\PP^2,$ $y'\in \PP^2\smallsetminus (S\cap V(y_1)).$
 It is linear in $A(x')A(y')^{-1}$ and is preserved under the products.
   Thus this condition is satisfied for the whole  algebra of matrices generated
 by (linear combinations and products of) matrices of type $A(x')A(y')^{-1}. $


Observe: the matrix  $A(y')^{-1}$ belongs to the algebra generated by $A(y').$ (By the Cayley-Hamilton Theorem.)
 Therefore the algebra generated by all the matrices of type $A(x'_o)A(y'_o)^{-1} $ coincides with the   algebra generated by $A_2,A_3,A_4.$


Finally   (a numerical computation): this algebra generated by $A_2,A_3,A_4$ is of dimension 36, i.e. it coincides with the
 vector space $Mat_{6\times 6}(\C).$
 Therefore the condition $g\cdot A(x')A(y')^{-1}=A(x')A(y')^{-1}\cdot g$ implies: $g$ is   a scalar matrix.

Finally, by scaling, we can assume $g=\Id.$
\end{proof}
\end{comment}



\begin{thebibliography}{CFFL24}
\bibliographystyle{plain}

\bibitem[BDG24]
{BDG}
Mikhail Borovoi, Bodgan A. Dina, and Willem A. de Graaf,
{\em Real non-degenerate two-step nilpotent Lie algebras of dimension eight},
Eur. J. Math. {\bf 10} (2024), no. 1, Paper No. 16, 27 pp.

\bibitem[Dem26]
{Demarche}
Cyril Demarche,
{\em A note on complex Lie algebras isomorphic to their conjugate},
arXiv:2604.06979.

\bibitem[Der19]
{Dere}
Jonas Deré,
{\em The structure of underlying Lie algebras},
Linear Algebra Appl. {\bf 581} (2019), 471--495.

\bibitem[GT99]
{GT}
L.~Yu.~ Galitski and Dmitry A.~Timashev.
{\em  On classification of metabelian Lie algebras},
J. Lie Theory {\bf 9} no. 1, 125--156 (1999).

\bibitem[GG20]
{GG20} Skip Garibaldi and Robert M.  Guralnick, {\em Generically free
representations, I: Large representations},
Algebra Number Theory {\bf 14} (6)) 1577--1611, 2020.

\bibitem[GL24]
{GL24}
Robert M. Guralnick and Ross Lawther,
{\em Generic stabilizers in actions of simple algebraic groups},
Mem. Amer. Math. Soc. {\bf 300} (2024), no. 1502, v+303 pp.


\bibitem[Lan02]
{Lang}
Serge Lang,
{\em Algebra},
Revised third edition,
Grad. Texts in Math., 211,
Springer-Verlag, New York, 2002. xvi+914 pp.

\end{thebibliography}


\end{document}


